\section{خلاصهٔ اجرایی}
این پروپوزال تأسیس یک آزمایشگاه پژوهشی را پیشنهاد می‌کند که سه حوزهٔ مرتبط با امنیت و سنجش کوانتومی را در یک زیرساخت مشترک پوشش می‌دهد:
\begin{enumerate}[leftmargin=*,itemsep=2pt]
\item \textbf{رمزنگاری پساکوانتومی \lr{(PQC)}:} الگوریتم‌های کلاسیک مقاوم در برابر رایانهٔ کوانتومی، همراه با ارزیابی امنیت پیاده‌سازی (حملات کانال جانبی و تزریق خطا).
\item \textbf{توزیع کلید کوانتومی \lr{(QKD)}:} ساخت یک ستاپ آزمایشی با متغیرهای پیوسته \lr{(CV-QKD)} روی فیبر نوری.
\item \textbf{سنجش و رادار کوانتومی:} آزمایش نوری روشن‌سازی کوانتومی \lr{(Quantum Illumination)} در مقیاس آزمایشگاهی.
\end{enumerate}

سه حوزه به‌ترتیب «سرمایه‌گذاری و پیچیدگی فزاینده» و مرحله‌ای اجرا می‌شوند و بین فازها دروازهٔ تصمیم قرار دارد؛ یعنی فاز بعدی فقط با تحقق معیارهای فاز قبلی و تأمین بودجه آغاز می‌شود.

\begin{longtable}{C{1cm}R{3.4cm}R{2.49cm}R{8cm}}
\toprule \rowcolor{hdrbg} فاز & موضوع & زمان پیشنهادی & خروجی اصلی \\ \midrule
\endfirsthead
\toprule \rowcolor{hdrbg} فاز & موضوع & زمان پیشنهادی & خروجی اصلی \\ \midrule
\endhead
$0$ & راه‌اندازی زیرساخت و تیم & ماه $1$ تا $3$ & فضا، ایمنی، تیم هسته، نقشهٔ تأمین، قرارداد تجهیزات فاز $1$. \\
$1$ & \lr{PQC} & ماه $3$ تا $12$ & بنچمارک و پیاده‌سازی \lr{ML-KEM/ML-DSA} روی سخت‌افزار هدف، آزمایشگاه حملات جانبی، آزمایش یکپارچه‌سازی \lr{TLS/PKI}. \\
$2$ & \lr{CV-QKD} & ماه $9$ تا $24$ & ستاپ تمام‌فیبری با مدولاسیون گاوسی، ثبت و پردازش آفلاین، کالیبراسیون نویز شات، گزارش کلید با پارامترهای امنیت. \\
$3$ & روشن‌سازی کوانتومی (نوری) & ماه $21$ تا $36$ & منبع جفت‌فوتون، اندازه‌گیری همبستگی در حضور نویز، مقایسه با مبنای کلاسیک. \\
\bottomrule\end{longtable}



